\documentclass[11pt]{article}
\usepackage[T1]{fontenc}
\usepackage[margin=1in]{geometry}
\usepackage{amsmath,amssymb,amsthm}
\usepackage[colorlinks=true,allcolors=blue]{hyperref}
\newtheorem{theorem}{Theorem}
\newtheorem{lemma}[theorem]{Lemma}
\title{Complex zeros and biased noise for the truth tables accepted by MCSP}
\author{Research manuscript}
\date{22 November 2026}
\begin{document}\maketitle
\begin{abstract}
Each truth table accepted by Boolean MCSP contributes one monomial, whose degree is its Hamming distance from a chosen center. Write $N=2^n$ for the table length and $s$ for the gate budget. Suppose $n=o(s)$ and $s=o(N/n)$. Define $K$ by $K\log_2(N/K)=s\log_2 s$, choosing the smaller solution. Around every center of complexity $o(s)$, the nearest zero is asymptotic to $-W(1/e)K/N$, and the rescaled zeros have the classical Szeg\H{o} limit distribution. For every fixed $0<\theta<1$, the minimum zero modulus among centers computable with $\theta s$ gates is $(W(1/e)+o(1))(1-\theta)K/N$. Exponentially many centers attain this asymptotic scale. Independent noise that flips about $K$ entries crosses the acceptance threshold at a center of complexity $o(s)$. If $K^3=o(N)$, the nearest zero at a uniformly chosen accepted center is farther away by a factor tending to infinity in probability. When $s\log s=o(\sqrt N)$ there is also a center of complexity exactly $s$ around which every zero approaches the unit circle. At a fixed fraction of the Shannon budget, a typical center on a fixed positive weight layer has three positive-axis points of zero accumulation. The limiting number of zeros in a small disk grows quadratically in its radius at the two noise transitions and linearly at one.
\end{abstract}
\section{Statement and statistical interpretation}
Put $N=2^n$. Circuits have unrestricted fanout, arbitrary binary Boolean gates of unit cost, and free constants and complemented wires. Let $\mathcal A=\{f:C(f)\le s\}$ count distinct truth tables. For $L\in\mathcal A$ put $P_L(z)=\sum_{f\in\mathcal A}z^{d(f,L)}$. If $L_p$ is formed by independently flipping the entries of $L$ with probability $p$, then for $0\le p<1$,
\[
 \Pr[L_p\in\mathcal A]=(1-p)^N P_L(p/(1-p)).
\]
Set $F=P_0$ and $a_k=[z^k]F$. All zero counts include multiplicity; $\log$ denotes the natural logarithm unless a base is indicated.
\begin{theorem}
Assume $n=o(s)$ and $s=o(N/n)$. Define $K$ as the smaller positive solution of $K\log_2(N/K)=s\log_2 s$. Let $L$ be any sequence of centers with $C(L)=o(s)$. Each nearest zero $z$ of $P_L$ satisfies
\[
                  \frac{Nz}{K}\longrightarrow-W(1/e),
\]
where $W$ is the positive inverse of $x\mapsto xe^x$. The normalized zero measure at this scale is given in Section~2. If $p=xK/N$ for fixed $x>0$, then
\[
 \lim \frac1K\log\Pr[L_p\in\mathcal A]=
 \begin{cases}0&x\le1,\\1+\log x-x&x>1.\end{cases}
\]
For $0<x<1$, the complementary probability has
$K^{-1}\log\Pr[L_p\notin\mathcal A]\to1+\log x-x$.
\end{theorem}
For polynomially bounded $s$, $K\sim s\log_2(s)/n$. The implicit scale matters for larger budgets. For example, if $s=N^{a+o(1)}$ with fixed $0<a<1$, then $K\sim as/(1-a)$.
\begin{theorem}
If also $s\log s=o(\sqrt N)$, some center $L_*$ of complexity exactly $s$ has every zero of $P_{L_*}$ at modulus $1+o(1)$, and its zeros approach the unit circle in Hausdorff distance. At $p=xK/N$ its acceptance probability is $(1+o(1))(1-p)^N$. Consequently its logarithmic rate is $-x$.
\end{theorem}
The proof of the noise estimates and the scale selection follows in Section~3. The second theorem expresses how a change of center can alter both roots and acceptance probabilities for the same circuit-size threshold.
\section{Transfer of the local Szeg\H{o} law}
Suppose \(K=K_N\to\infty\), \(K/N\to0\), and

\par\noindent (i) For every fixed \(\delta>0\), eventually \(a_k={N\choose k}\) for \(0\le k\le(1-\delta)K\).
\par\noindent (ii) \(\log F(1)\le (1+o(1))K\log(N/K)\).

Set \(G_N(w)=F(Kw/N)\), \(u_N=K^{-1}\log|G_N|\), \(r_*=W(e^{-1})\), and let
\[
 D=\{0\}\cup\{0<|w|<1:\ 1+\log|w|<\Re w\}.
\]
For compact \(E\subset D\), \(G_N(w)/(1+Kw/N)^N=1+O_E(e^{-c_E K})\) uniformly on \(E\). In \(L^1_{\rm loc}(\mathbb C)\), \(u_N\to U\), with \(U=\Re w\) on \(D\) and \(U=1+\log|w|\) outside. The local counting measures \(K^{-1}\sum_{G_N(w)=0}\delta_w\) converge vaguely to \(\Delta U/(2\pi)\), a probability measure on \(\partial D\) with arclength density \(|1/w-1|/(2\pi)\). Consequently the least zero modulus of \(F\) is \((r_*+o(1))K/N\). Every fixed neighborhood of a regular boundary arc contains order \(K\) local zeros. Every closest zero, scaled by \(N/K\), tends to \(-r_*\). These statements allow \(o(K)\) outliers on compact subsets of the exterior.

\textbf{Interior estimate.} Fix \(r=|w|<1\), \(m=\lfloor(1-\delta)K\rfloor\), with \(\delta\) small and \(r<1-\delta\). The difference from the full binomial is bounded by
\[
 \left|(1+Kw/N)^N-G_N(w)\right|
 \le \sum_{k>m}{N\choose k}(Kr/N)^k.
\]
Above \(m\) the ratio of successive terms is bounded away from one. Uniform Stirling gives limiting upper exponential rate \( (1-\delta)(1+\log(r/(1-\delta)))\) on scale \(K\), tending to \(1+\log r\). The full binomial has rate \(\Re w\), locally uniformly since \(K=o(N)\). This proves exponential approximation throughout compact subsets of \(D\) (bound on a small disk separately at zero). In \(r<1\) the same argument gives the upper rate \(\max(\Re w,1+\log r)\).

\textbf{Exterior saturation.} On compact disks \(|w|\le R\), the terms with \(k>(1+\delta)K\) are bounded by \(F(1)(KR/N)^{(1+\delta)K}\); their limiting upper logarithmic rate is minus infinity. If \(r>1\), choose \(0<\delta<r-1\). Terms of the remaining binomial majorant increase up to \((1+\delta)K\), so Stirling and \(\delta\downarrow0\) give upper rate \(1+\log r\). At \(r=1\) use nearby radial bounds. At a positive real \(x>1\), the complete single layer \(\lfloor(1-\delta)K\rfloor\) attains this rate as \(\delta\downarrow0\). The locally upper bounded subharmonic functions \(u_N\) are precompact in \(L^1_{\rm loc}\), collapse to minus infinity being excluded by the interior estimate. For each subsequential limit \(u\), its upper semicontinuous representative lies below the exterior harmonic upper bound by the submean inequality; Hartogs' lemma yields \(u(x)\ge\limsup u_N(x)\) at the fixed positive \(x\). The strong maximum principle forces equality on the connected exterior. Indeed \(1+\log r-r\cos\theta\) increases strictly for \(0<r<1\), so \(D\) is star-shaped with a Jordan boundary of zero area. This proves \(L^1_{\rm loc}\) convergence, and the distributional Laplacian gives the zero measure. Its arc density follows from the gradient jump; flux across a large circle gives mass one. On \(\partial D\) the minimum radius is attained only at \(-r_*\); compact zero exclusion inside and positive limit mass near that point give the nearest-zero assertions.
\section{What the circuit-size bound gives}
A trimmed, rooted circuit with $s$ gates can be encoded by depth-first traversal from its output. At most $s+1$ inputs of gates are pointers to an already introduced gate or a raw input, because $s-1$ of the at most $2s$ gate inputs introduce a new gate. Pointer types, traversal and gate labels take $O(s)$ extra bits. Including the zero-gate circuits and all smaller sizes, this gives
\[
 \log_2 |\mathcal A|\le (s+1)\log_2(s+n+2)+O(s+\log n).
\]
Lupanov's fixed-weight synthesis theorem \cite{lup,kors} bounds the maximum cost among tables of weight $\ell$ by $(1+o(1))H/\log_2 H$, where $H=\log_2\binom{N}{\ell}$ and $\min(\ell,N-\ell)/\log n\to\infty$. Synthesis on the standard binary basis also gives an upper bound in our full binary basis.
Write $S=s\log_2 s$. Since $S=o(N)$, the increasing branch of $k\mapsto k\log_2(N/k)$ gives $K/N\to0$ and $\log(N/K)\to\infty$. Also $K\ge S/n\gg\log n$, because $s/n\to\infty$. For each fixed $t>0$, Stirling's formula yields
\[
 \log_2\binom{N}{\lfloor tK\rfloor}=(t+o(1))S,
 \qquad \log_2\binom{2N}{\lfloor tK\rfloor}=(t+o(1))S .
\]
The extra $O(K)$ from doubling the domain is negligible. Put $\ell=\lfloor(1-\delta)K\rfloor$ for fixed $0<\delta<1$. Chashkin's small-on-sets synthesis bound \cite[Theorem~6.1]{chash} applies on the full domain to every perturbation of weight at most $\ell$. Its cost is at most $(1-\delta+o(1))s$, since the entropy is $(1-\delta+o(1))S$, $\log_2 S\sim\log_2 s$, and the $O(n)$ overhead is $o(s)$. Alternatively, append an address bit, take any table of weight $k\le\ell$ in its zero half and place $\ell-k$ ones in its other half. Lupanov synthesis at weight $\ell$ on the doubled domain has the same budget. Fixing the extra bit recovers the original table. XOR with $L$ costs at most $C(L)+1=o(s)$ additional gates. Thus every translated layer through $(1-\delta)K$ is complete. The rooted count gives $\log|\mathcal A|\le(1+o(1))K\log(N/K)$ for every center. Apply Section~2 with $F=P_L$.

For noise $p=xK/N$, the number of flips is binomial with mean $xK$. At the single layer $k=\lfloor tK\rfloor$ its probability has log rate $-x+t+t\log(x/t)$ after division by $K$, by Stirling. When $x<1$, rejection is bounded above by the binomial tail above $(1-\delta)K$, taking $0<\delta<1-x$. At layer $(1+\delta)K$, the fraction of accepted tables is at most $2^{-(\delta+o(1))S}$ by the rooted count. This layer gives the lower bound for rejection. Send $\delta$ to zero. For $x>1$, multiply the positive real-axis rate from Section~2 by the binomial normalization $(1-p)^N$, using $p/(1-p)=(x+o(1))K/N$. A positive-axis squeeze between fixed nearby values justifies this substitution and gives the endpoint $x=1$ by continuity.

For $s\sim cN/n$ with fixed $0<c<1$, let $\alpha\in(0,1/2)$ solve $H_2(\alpha)=c$, with $H_2$ the base-two binary entropy. For fixed $\delta>0$, synthesize every weight between $\epsilon N$ and $(\alpha-\delta)N$ directly by Lupanov, leaving a fixed fractional gate margin. For all smaller weights, pad as above to weight $\lceil\epsilon N\rceil$ on $n+1$ inputs, choosing $\epsilon$ sufficiently small that $2H_2(\epsilon/2)<c$. Adding a center of cost $o(s)$ by XOR fits the margin. Complementation gives both complete tails. The rooted count supplies $\log P_L(1)\le(c\log2+o(1))N$.
\section{The fixed-fraction zero measure}
Fix \(0<\alpha<1/2\), \(h(t)=-t\log t-(1-t)\log(1-t)\), \(r_\alpha=\alpha/(1-\alpha)\). Assume \(\log F(1)\le Nh(\alpha)+o(N)\), and every weight \(k\le(\alpha-\delta)N\) or \(k\ge(1-\alpha+\delta)N\) is complete for every fixed \(\delta>0\) eventually. Define
\[
 B_\alpha=\{0\}\cup\{0<|z|<r_\alpha:\ h(\alpha)+\alpha\log|z|<\log|1+z|\}.
\]
In \(|z|<1\), the \(L^1_{\rm loc}\) limit of \(N^{-1}\log|F(z)|\) equals \(V(z)=\log|1+z|\) in \(B_\alpha\) and \(h(\alpha)+\alpha\log|z|\) elsewhere. In \(|z|>1\) it is \(\log|z|+V(1/z)\). These formulas extend continuously to the common value \(h(\alpha)\) on the unit circle. The global normalized zero measure converges weakly to three components:

\par\noindent$\bullet$\ arclength density \((2\pi)^{-1}|\alpha/z-1/(1+z)|\) on \(\partial B_\alpha\), with mass \(\alpha\);
\par\noindent$\bullet$\ its pushforward by \(z\mapsto1/z\), also of mass \(\alpha\);
\par\noindent$\bullet$\ uniform unit-circle measure of mass \(1-2\alpha\).

The nearest modulus tends to the unique \(t\in(0,r_\alpha)\) with \(\log(1-t)=h(\alpha)+\alpha\log t\), and every closest zero tends to \(-t\).

\textbf{Proof.} For \(r=|z|<r_\alpha\) the binomial tail after \(m=\lfloor(\alpha-\delta)N\rfloor\) is bounded by \(O({N\choose m}r^m)\) when \(\delta\) is sufficiently small; its limiting upper rate tends to \(h(\alpha)+\alpha\log r\) as \(\delta\downarrow0\). This gives exponential approximation by the full binomial in the lobe and the stated upper bound outside. For \(r_\alpha<r<1\), split at \(m=\lfloor\alpha N\rfloor\). The lower part is bounded by \((N+1){N\choose m}r^m\) since its binomial majorant increases, and the upper part by \(F(1)r^m\). Both have limiting upper rate \(h(\alpha)+\alpha\log r\). At a fixed positive \(r\) in this annulus, a single complete layer \((\alpha-\delta)N\) attains that rate in the limit \(\delta\downarrow0\). The exterior is connected within the unit disk: for \(r<r_\alpha\),
\[
 \frac{d}{d\log r}(\log|1+re^{i\theta}|-h(\alpha)-\alpha\log r)
 =\Re\frac{z}{1+z}-\alpha\le\frac r{1+r}-\alpha<0 .
\]
At \(r=r_\alpha\) the inequality defining the lobe becomes equality only on the positive axis, and fails elsewhere. The same subharmonic maximum principle now proves the disk result. Apply it to the reciprocal polynomial to cover \(|z|>1\). Exact complement symmetry is unnecessary; high-weight completeness suffices. Global convergence follows by local boundedness and identification off area-zero interfaces. The flux at radii \(r_\alpha<r<1\) is \(2\pi\alpha\); the radial slopes across the unit circle jump from \(\alpha\) to \(1-\alpha\). These give the masses above. As they sum to one, vague convergence implies tightness and weak convergence. The least boundary radius occurs on the negative real axis by minimizing \(|1+z|\) at each radius.

For comparison, at every $s=o(N/n)$ and every accepted center $L$, the unscaled zero measure, divided by $N$, tends weakly to uniform measure on the unit circle. Indeed $P_L$ has leading and constant coefficients one by complement symmetry and $\log P_L(1)=o(N)$ by circuit counting. Jensen's formula and reciprocity bound the number of zeros outside $e^{-\eta}<|z|<e^\eta$ by $2\log P_L(1)/\eta$. The Erd\H{o}s--Tur\'an angular discrepancy bound \cite{et} completes the proof. The $K=o(N)$ zeros detected near a low-complexity center are visible in the finer scale of Section~2.

\section{A local zero law just below twice the address length}
Here the gate budget is linear in the address length, outside the synthesis regime of Theorem~1. The center is the zero Boolean table. Let $s=2n-v$ with $v=o(n)$ and $\liminf\log v/\log n\ge1/2$. Define $J=\log_2(n/v)$. Then $J\to\infty$ and $J\le(1/2+o(1))\log_2 n$.

\begin{theorem}\label{th:linear}
At these budgets the zero-center polynomial $P_0(Jw/N)$ has the local limit potential, zero measure and nearest-zero location of Section~2, with $J$ in place of $K$. The logarithmic acceptance and rejection rates of Theorem~1 hold at flip probability $xJ/N$ on scale $J$. The conclusions hold for full binary gates with free complemented wires and also for binary AND, OR and counted unary NOT.
\end{theorem}
\begin{proof}
For a nonempty set $S$ of $k$ addresses write $Z$ for the number of raw coordinates constant across its $k$ members. We use the sparse-support estimate
\[
       2n-Z-k-\log_2 k-1 \le C({\bf1}_S)
                    \le 2n-Z+O(k2^k+1).                         \tag{4}
\]
For the lower bound, the companion sparse gate-elimination lemma in \emph{Unrestricted Boolean circuit lower bounds for linear-threshold MCSP}, Section~1, gives $C\ge n+b-r-\nu_2(k)-1$ for full binary gates, where $b=n-Z$ and $r=\dim\operatorname{aff}(S)\le k-1$. This implies the displayed lower bound for both bases. For the upper bound group the raw coordinates by their columns of values across the $k$ support points. A nonconstant column group of length $m$ needs at most $2m$ gates to test whether all its bits are one and whether all its bits are zero. One AND chain, one OR chain and one NOT suffice. A constant group needs only its required test, at cost at most $m$. For each support point conjoin the correct group tests, and take the OR of these $k$ terms. There are at most $2^k$ groups, giving (4). Red'kin's work on functions with few ones also groups variables according to their behavior on the support \cite{red}.

Choose $k$ independent uniform addresses and then condition on distinctness. For $k=O(\log n)$ conditioning costs a factor $1+o(1)$ in any probability upper bound. Before conditioning, $Z$ is binomial with mean $\mu_k=2n\,2^{-k}$. For any fixed $0<\eta<1$, uniformly at $1\le k\le(1-\eta)J$ one has $\mu_k/v\to\infty$ and $\mu_k/(k2^k+1)\to\infty$. Indeed $2^{2k}\le n^{1-\eta+o(1)}$ there. Thus (4) and a binomial lower-tail bound show that the rejected fraction of each such shell is at most $\exp(-\Omega(v))$. Conversely for fixed $B>1+\eta$ and $(1+\eta)J\le k\le BJ$, $\mu_k/v\to0$ and $k=o(v)$. The lower bound in (4) forces $Z\ge v-O(k+\log k)$ for acceptance. The binomial upper tail bounds the accepted fraction by $\exp(-\Omega(v))$, uniformly on this range.

Here the two shell-fraction estimates replace exact early shells and the global census. On each fixed disk $|w|\le b$, the sum of the errors in the early binomial coefficients, multiplied by $(bJ/N)^k$, is at most $\exp(-\Omega(v)+bJ)$. The same bound holds for the contribution from $(1+\eta)J$ to $BJ$. Since $v/J\to\infty$ these errors have logarithmic rate minus infinity on scale $J$. For $k>BJ$ use the full binomial majorant. Its tail has upper rate at most $B(1+\log(b/B))$ when $B>b$, by the geometric ratio and Stirling's formula. This tends to minus infinity as $B\to\infty$. The interior comparison and exterior saturation proof of Section~2 now apply with the same potential. The real-axis lower bound uses a layer with accepted fraction $1-o(1)$.

For rejection below the noise threshold, bound failures among the early shells by $\exp(-\Omega(v))$ and the remaining probability by the binomial upper tail. A layer just above $(1+\eta)J$ is rejected with probability $1-o(1)$ conditional on its size. These upper and lower bounds give the rate as $\eta\downarrow0$. The acceptance rates come from the positive-axis potential.
\end{proof}

There is also an all-center statement far below this regime. If $s+1\le n$, the algebraic degree over $\mathbb F_2$ of every $F\in\mathcal A$ is at most $s+1$, because $F$ depends on at most $s+1$ address variables. Therefore the difference of two distinct easy tables has weight at least $2^{n-s-1}$, by the minimum-distance bound for degree-$(s+1)$ multilinear polynomials. One direct proof restricts a variable of a highest-degree monomial. The XOR of the two restrictions has degree one smaller, and its weight is at most the sum of their weights. Induction on the degree, starting with a nonzero constant, gives the bound. Where output complementation is free, for every accepted center all zeros consequently have
\[
                |\log|z||\le 2^{-(n-s-1)}\log |\mathcal A|.
\]
Inside the indicated inner radius the other terms have sum strictly less than the constant term, and reciprocity gives the outer radius. Thus whenever $2^{n-s}\gg (s+1)\log(n+s)$ all centers have every zero at modulus $1+o(1)$.

\section{The last partially filled layer and the nearest zero}
Here \(k=k(n)\to\infty\), \(k^2 8^k=o(n)\), \(N=2^n\), and \(C_k\) denotes the minimum gate count of a uniformly chosen truth table of weight \(k\). Put
\[
 \mu=n2^{1-k},\qquad \sigma^2=\mu(1-2^{1-k}).
\]
For these assertions it suffices to have the sparse-support bounds \(2n-Z-O(k+\log k)\le C\le2n-Z+O(k2^k)\), where \(Z\) counts constant raw columns of the \(k\) support addresses. The full binary basis and counted De Morgan basis obey these inequalities by the sparse-support lower bound and repeated-column construction used above. Red'kin's \href{https://www.mathnet.ru/eng/dm172}{2004 theorem} analyzes the Finikov synthesis method over the binary basis with XOR and XNOR excluded.

For independent random rows, \(Z\) is binomial with mean \(\mu\). Conditioning them to be distinct has total-variation cost at most \(\binom{k}{2}/N\). The binomial Berry--Esseen inequality, applied to \(-Z+\mu\), and the two deterministic shifts of size \(O(k2^k)\) therefore give
\[
 \sup_x\left|\Pr\!\left[\frac{C_k-(2n-\mu)}{\sigma}\le x\right]
        -\Phi_{\rm G}(x)\right|
 \ \le O\!\left(\frac{k2^k+1}{\sigma}+\frac{k^2}{N}\right)=o(1),
\]
where \(\Phi_{\rm G}\) is the standard normal distribution function. Translating the Gaussian argument costs at most the shift times its bounded density.

Take \(s=\lfloor2n-\mu+x_n\sigma\rfloor\), with \(|x_n|=o(\sqrt\mu)\), and let \(\lambda_k=\Pr[C_k\le s]\) be the \emph{actual} accepted fraction on layer \(k\). Write \(r=W(1/e)\), \(b=r/(1+r)\), \(\lambda_*=1/(1+r)\) and \(a=r/(1+r)^2\). Let \(R_N\) be the least modulus of a zero of \(P_0\) at this budget.
\begin{theorem}
If \(\sqrt{k}\,|\lambda_k-\lambda_*|\to\infty\), then
\[
 NR_N=rk+b\left[\tfrac12\log(2\pi k)-\log|\lambda_k-\lambda_*|\right]+o(1).
\]
Every nearest zero \(\zeta\) satisfies \(\sqrt{k}\,|\arg(-\zeta)|\to0\) in this case; \(\arg\) is the principal argument. If \(\sqrt{k}(\lambda_k-\lambda_*)\to h\in\mathbb R\), put
\[
 B_h(t)=\frac{r t^2}{2}-\frac12\log(h^2+a^2t^2),\qquad
 t_h=\sqrt{\max(1/r-h^2/a^2,0)} ,
\]
interpreting a zero logarithm argument as \(B_h(t)=+\infty\). Then
\[
 NR_N=rk+b\left[\log k+\tfrac12\log(2\pi)
                       +B_h(t_h)\right]+o(1),\qquad
       \sqrt{k}\,|\arg(-\zeta)|\longrightarrow t_h .
\]
For instance the Gaussian approximation can be used in place of \(\lambda_k\) at the critical precision if \(k^3 8^k=o(n)\).
\end{theorem}

\begin{proof}
Let \(E_j(u)=\sum_{i=0}^j u^i/i!\) and \(G_k(w)=P_0(kw/N)\). On each fixed disk \(|w|\le R\), for every fixed \(A>0\),
\[
  G_k(w)=E_{k-1}(kw)+\lambda_k(kw)^k/k!+O_{A,R}(e^{-Ak}).
\]
Indeed the constant-column mean is at least \(2\mu\) below layer \(k\) and at most \(\mu/2\) above it. In any range \(1\le j\le Dk\), for fixed \(D\), the lower-layer upper synthesis error is \(O(k2^k)=o(\mu)\) and the upper-layer gate lower-bound error is \(O_D(k+\log k)=o(\mu)\). Chernoff bounds and distinctness conditioning thus make the incorrect fraction \(O(e^{-\Omega(\mu)}+(Dk)^2/N)\) outside layer \(k\). Layer zero is complete. The binomial-to-Poisson coefficient error through \(Dk\) is \(O_D(k^2/N)\) relative to the corresponding exponential-series coefficient before thinning. These errors are multiplied by at most \(e^{Rk}\). For the farther terms use the full exponential-series majorant and take \(D\) large enough that \(D(1+\log(R/D))<-A-1\).

On closed annuli \(0<\eta\le|w|\le R<1\) write exactly
\[
 G_k(w)=e^{kw}-\frac{(kw)^k}{k!}A_k(w),\qquad
 A_k(w)=\frac1{1-w}-\lambda_k+O_{\eta,R}(1/k).
\]
The same error estimate holds after one or two differentiations on smaller annuli. Indeed the normalized exponential tail is
\(\sum_{j\ge0}w^j\prod_{i=1}^j(1+i/k)^{-1}\), whose difference from \((1-w)^{-1}\) is bounded by \(\sum_{j\ge0}R^jj(j+1)/(2k)\). The preceding arbitrarily small exponential error remains negligible on the annulus after normalization. Cauchy's inequality gives the derivative bounds. The nonvanishing of \(e^{kw}\) and the same tail estimate exclude every disk \(|w|\le r-\eta\) from the zeros for large \(k\): there \(|w|e^{1+|w|}<1\) uniformly.

Write \(w=-\rho e^{i\theta}\) near \(-r\), \(d=k(\rho-r)\), and \(t=\sqrt k\,|\theta|\). For a zero, Stirling's formula gives the exact modulus equation up to a uniform \(O(k^{-1})\),
\[
 k(\Re w-1-\log\rho)=\log|A_k(w)|-\tfrac12\log(2\pi k)+O(k^{-1}).
\]
For \(|d|=O(\log k)\) and \(t=O(\sqrt{\log k})\) it becomes
\[
 b^{-1}d=\tfrac12\log(2\pi k)-\log|A_k(w)|+\tfrac12 rt^2+o(1).
 \tag{*}
\]
The error is uniform because \(k\theta^4+ d^2/k+|d|\theta^2=o(1)\) and \(-r-1=\log r\).

We first verify that the modulus equation really locates zeros to \(o(1)\) in \(d\). In the noncritical regime take \(\theta=0\); in the critical regime take \(\theta=t_h/\sqrt k\). In the latter case \(h^2+a^2t_h^2>0\). Uniformly through \(|d|\le C\log k\), the corresponding tail magnitude is asymptotic respectively to \(|\lambda_*-\lambda_k|\) or to \(k^{-1/2}\sqrt{h^2+a^2t_h^2}\): Taylor expansion of \((1-w)^{-1}\) gives the imaginary linear term \(-ia\theta\), with error \(O(\theta^2+|d|/k+1/k)\). Consequently the exact modulus equation has a solution \(w_0\) at the asserted value of \(d\) up to \(o(1)\), by the intermediate value theorem. More explicitly, insert into (*) the asserted expression for \(d\) plus and minus any fixed positive tolerance; the two residuals have opposite signs.

Choose analytic logarithms near \(w_0\) and write
\[
 T_k(w)=kw-k\log w-\log A_k(w)+\log k!-k\log k.
\]
The exact modulus equation is \(\Re T_k(w_0)=0\). A zero results whenever \(T_k(w)\in2\pi i\mathbb Z\). At \(w_0\), \(T'_k(w_0)=k(1+1/r)+O(\sqrt k+\log k)\), since \(|A_k(w_0)|\gtrsim k^{-1/2}\). Throughout a disk of radius \(C'/k\) for any fixed \(C'\) we have \(T''_k=O(k)\). Take \(2\pi i m\) nearest to \(T_k(w_0)\). The linearized root differs from \(w_0\) by \((2\pi i m-T_k(w_0))/T'_k(w_0)\), of magnitude \(O(1/k)\) and real part \(O(k^{-3/2})\). Taylor's theorem and Rouch\'e's theorem on a circle of radius \(C''/k^2\) about this linearized root, with sufficiently large \(C''\), give an exact root there. Its radial displacement from \(w_0\) is \(o(1/k)\), since \(\theta=O(k^{-1/2})\). This proves the stated expressions as upper bounds for \(NR_N\).

For a nearest zero these upper bounds give \(\rho\le r+O(\log k/k)\). The small-disk exclusion and the modulus equation, using \(\Re w\ge-\rho\) and boundedness of \(A_k\), now give \(d\ge-O(1)\) and \(k(1-\cos\theta)=O(\log k)\). Thus (*) applies to every nearest zero. Put \(H_k=\sqrt k\,(\lambda_k-\lambda_*)\). In this range
\[
 \sqrt k\,|A_k(w)|\le |H_k|+O(1+t).
\]
If \(|H_k|\to\infty\), comparison with the obtained upper bound shows \(rt^2/2\le\log(1+O(1+t)/|H_k|)+o(1)\). This first bounds \(t\) and then forces \(t\to0\), proving the matching radius by (*). For bounded \(H_k\), the same comparison yields \(rt^2/2\le\log(1+t)+O(1)\), so \(t\) is bounded again. A subsequential limit \(t'\) cannot satisfy \(h=t'=0\), since \(\sqrt k|A_k|\) would tend to zero in (*) in contradiction with the upper bound. Otherwise \(\sqrt k|A_k|\to\sqrt{h^2+a^2t'^2}\). Equation (*) now gives the correction \(B_h(t')\). Differentiation shows its unique minimum on \(t\ge0\) occurs at \(t_h\), also when \(h^2=a^2/r\). The upper bound forces this minimum for every subsequential limit, as claimed.
\end{proof}

For ordinary exponential sections, logarithmic corrections and the locations of individual zeros relative to the Szeg\H{o} curve have been studied by \href{https://doi.org/10.1216/rmjm/1181072998}{Carpenter, Varga and Waldvogel} and \href{https://www.math.kent.edu/~varga/pub/paper_178.pdf}{Varga and Carpenter}. Here the last coefficient can cancel the leading tail at the negative extremum, producing the critical scaling in the theorem.

\section{How the budget spent on the center changes the first zero}
Put $T=s\log_2 s$ and $r_*=W(e^{-1})$. In this section $n=o(s)$ and $s=o(N/n)$.
For each fixed $0<\theta<1$, define $J$ on the smaller branch by
\[
                J\log_2(N/J)=(1-\theta)T .
\]
Since $\log(N/K)\to\infty$, evaluation at fixed multiples of $K$ gives $J\sim(1-\theta)K$.
\begin{theorem}\label{th:profile}
Write $R_L$ for the least zero modulus of the distance enumerator centered at $L$. Then
\[
 \min_{C(L)\le\lfloor\theta s\rfloor}R_L
      =(r_*+o(1))\frac{J}{N}.
\]
For every sequence with $\limsup C(L)/s\le\theta$, the rescaled polynomial $P_L(Jw/N)$ has no zeros on any prescribed compact subset of $D$ for all sufficiently large $n$. Moreover, at least $2^{(\theta-o(1))T}$ distinct centers of complexity between $(\theta-o(1))s$ and $\lfloor\theta s\rfloor$ attain the local zero measure and the closest-zero location of Section~2 with $J$ in place of $K$. They also have all the noise rates in Theorem~1 with this substitution.
\end{theorem}
\begin{proof}
For every center in the universal statement and fixed $\delta>0$, each layer up to $(1-\delta)J$ is complete. Indeed the synthesis bound of Section~3 builds any such error pattern for at most $((1-\delta)(1-\theta)+o(1))s$ gates. XOR with the center fits the budget. The interior estimate in Section~2 only used these complete layers and the binomial majorant. This gives the zero exclusion and a uniform lower bound for the minimum by passage to an arbitrary minimizing sequence.


A stronger probabilistic statement supplies the upper bound. Fix any baseline $G$ of complexity $o(s)$. Take $H$ uniformly among sets of $m=(\theta+o(1))K$ positions and set $L=G\oplus {\bf1}_H$. With probability tending to one these centers obey all the asserted zero and noise laws, with convergence in probability for the numerical limits. The weight-$m$ synthesis estimate gives $C(L)\le(\theta+o(1))s$ uniformly. Its error and $C(G)$ can be absorbed by choosing $m=(\theta-\epsilon_n)K+O(1)$ for a sufficiently slowly decreasing $\epsilon_n>0$. This gives the literal gate cap in the theorem. There are $\binom Nm=2^{(\theta+o(1))T}$ distinct centers. The circuit census at budget $(\theta-\eta_n)s$ excludes a vanishing fraction for some $\eta_n\downarrow0$, chosen slowly enough to dominate the errors in both counts.

To prove the probabilistic claim write $a_d(L)=[z^d]P_L(z)$. Independently choose a uniform $d$-set $E$. For fixed $\delta>0$ and $(1+\delta)J\le d\le 2K$, consider the shell fraction
\[
 \mathbb E_L\frac{a_d(L)}{\binom Nd}
      =\Pr[G\oplus{\bf1}_{H\triangle E}\in\mathcal A].
\]
Here $J/K\to1-\theta$. Write $\gamma=\delta(1-\theta)/8$, shrinking $\delta$ if necessary. The chance that $|H\cap E|\ge\lceil\gamma K\rceil$ is at most $\binom{m}{\lceil\gamma K\rceil}(d/N)^{\lceil\gamma K\rceil}=2^{-\Omega_{\theta,\delta}(T)}$, with the value zero when $m$ is too small. On all other intersection sizes $|H\triangle E|\ge(1+2\gamma)K$ for large $n$. Conditional on its size, this symmetric difference is uniform on its layer. Its size is also $O(K)=o(N)$, so comparison of the rooted census with the layer size shows the same exponential bound for acceptance. These estimates are uniform in $d$ and $G$. By Markov's inequality and a union bound over the at most $2K+1$ layers, outside a vanishing fraction of $H$ all their accepted fractions are at most $2^{-c_{\theta,\delta}T}$ for a positive constant $c_{\theta,\delta}$.

Consequently on any disk $|w|\le b$, $b>0$ fixed, the contribution to $P_L(Jw/N)$ above degree $(1+\delta)J$ is bounded by
\[
 2^{-c_{\theta,\delta}T}(1+bJ/N)^N+
                 |\mathcal A|(bJ/N)^{2K}.
\]
The logarithm of each summand divided by $J$ tends to minus infinity. The first term uses $T/K\to\infty$, the second uses $K\log_2(N/K)=T$. Choose $\delta$ decreasing to zero sufficiently slowly by a diagonal argument to obtain one set of $1-o(1)$ of all centers on which every fixed margin works. The complete early layers and the exterior proof of Section~2 now give the potential, zero measure, and nearest-zero location. The positive-axis values give the acceptance rates. For subcritical rejection, the accepted fraction of layer $\lceil(1+\delta)J\rceil$ tends to zero, as just proved. The single-layer argument from Section~3 completes the proof.
\end{proof}

The theorem gives an extremal law over all centers below a prescribed circuit budget. It also supplies many centers with that law. The next section concerns a uniform center among all distinct accepted tables, whose local law can be different.

\section{Noise about a typical center of fixed positive weight}

Suppose \(sn/N\to c\in(0,1)\). Write \(H(x)=-x\log x-(1-x)\log(1-x)\) and let \(\alpha\in(0,1/2)\) solve \(H(\alpha)=c\log2\). Fix \(0<h<\alpha\) and any deterministic baseline \(B=B_N\) of cost \(o(s)\), including the zero table. Choose \(L=B\oplus U\), where \(U\) is uniform on a layer of integer weight \(h_NN\) with \(h_N\to h\). Every such center belongs to \(\mathcal A\) for all sufficiently large \(N\). Define
\[
 \Lambda_{h,p}(v)=(1-h)\log(1-p+pe^v)+h\log(p+(1-p)e^v),
 \qquad I_{h,p}(q)=\sup_v\{vq-\Lambda_{h,p}(v)\}.
\]
The critical noise and the corresponding positive argument are
\[
 p_0=\frac{\alpha-h}{1-2h}\in(0,1/2),\qquad r_0=\frac{p_0}{1-p_0}.
\]

\begin{theorem}
For \(0<p\le 1/2\), let \(A_L(p)=\Pr[L_p\in\mathcal A\mid L]\). In probability over the chosen center,
\[
 \frac1N\log A_L(p)\longrightarrow
 \begin{cases}0&p\le p_0,\\-I_{h,p}(\alpha)&p>p_0.\end{cases}
\]
For \(0<p\le p_0\), also \(N^{-1}\log(1-A_L(p))\to-I_{h,p}(\alpha)\) in probability. For each fixed \(p\) and positive error tolerance, the exceptional fraction is at most \(\exp(-\Omega(N))\), with a constant depending on those parameters. Complement symmetry gives \(A_L(1-p)=A_L(p)\).

Every fixed open disk about any of \(r_0,1,1/r_0\) contains at least \(\gamma N\) zeros of \(P_L\) with probability tending to one, for some \(\gamma>0\) depending on the disk and the fixed parameters. For \emph{every} center in the indicated layer there are eventually no zeros on compact subsets of
\[
 \left\{z:\ p=\frac{|z|}{1+|z|}<p_0,\quad
       \log\frac{1+|z|}{|1+z|}< I_{h,p}(\alpha)\right\}
\]
or its reciprocal image, interpreting the condition at zero by continuity.
\end{theorem}

\begin{proof}
Set \(\mathcal B=B\oplus\mathcal A\). The census gives \(|\mathcal B|\le\exp(NH(\alpha)+o(N))\). For every fixed \(\delta>0\), \(\mathcal B\) contains all layers of weight at most \((\alpha-\delta)N\), by synthesis with slack and combination with \(B\). Conditional on \(U\), the weight \(W\) of the perturbed \(U\) is a sum of two independent binomials of parameters \(((1-h_N)N,p)\) and \((h_NN,1-p)\). Its rate function is \(I_{h,p}\), strictly convex with minimum zero at \(\mu=h+(1-2h)p\). More explicitly, for \(qN\) integral with \(q\) tending to an interior point, the logarithmic probability of \(W=qN\), divided by \(N\), tends to \(-I_{h,p}(q)\) evaluated at that limit. This follows by maximizing the two binomial point probabilities subject to their counts summing to \(qN\), using Stirling's formula and a polynomial number of count pairs. The maximum is obtained by a common exponential tilt \(v\) with \(\Lambda'_{h,p}(v)=q\). The Chernoff bound gives the corresponding uniform upper rate, including the endpoints by continuity. These estimates hold for every \(U\) in the specified layer.

After averaging over \(U\), conditional on \(W\) the perturbed table is uniform in its layer by permutation invariance. The exponential upper rate for the mean acceptance is consequently the negative of
\[
 \min_{0\le q\le1}\{I_{h,p}(q)+(H(q)-H(\alpha))_+\}.
\]
For \(p>p_0\) this minimum is \(I_{h,p}(\alpha)\). Indeed \(\alpha<\mu\le1/2\). On \([\alpha,1-\alpha]\) the derivative of \(I_{h,p}(q)+H(q)\) is \(v-\log(q/(1-q))\), which is nonnegative throughout this interval. To check the sign for \(p<1/2\), put \(x=e^v\), \(a=p/(1-p)\). The sign of \(\Lambda'_{h,p}(v)-x/(1+x)\) is that of \(h-(1-h)a-(1-h-ha)x\). If this linear expression has a positive root, the associated crossing occurs at \(q=(h-p)/(1-2p)<h\); above it the expression is negative. For \(p=1/2\) the derivative is identically zero. On the two exterior intervals convexity and \(\alpha<\mu\le1/2\) put the minima at their inner endpoints; the middle-interval comparison orders those endpoints. When \(p\le p_0\), the minimum is zero at \(\mu\).

The complete layers supply the matching lower rate, for every center, by taking weights tending to \(\min(\mu,\alpha)\) from below. Markov's inequality turns the averaged upper rate into the claimed exponential exceptional-set bound. For \(p<p_0\), rejection is bounded above for every center by \(\Pr[W>(\alpha-\delta)N]\), which gives upper rate \(-I_{h,p}(\alpha)\) after \(\delta\downarrow0\); at \(p=p_0\) the upper rate zero is immediate. For the reverse estimate choose \(q=\alpha+\delta<1/2\). Conditional on \(W=\lfloor qN\rfloor\), the mean accepted fraction is exponentially small by the census. For all but an exponentially small fraction of centers the conditional rejected fraction exceeds \(1/2\). The point probability of this layer supplies the lower rate, taking \(\delta\) small for the specified tolerance.

For \(r>0\), let \(\Phi(r)\) be the limiting \(N^{-1}\log P_L(r)\), obtained by adding \(\log(1+r)\) to the acceptance exponent at \(p=r/(1+r)\) and using the stated symmetry on the other half. The convergence is locally uniform in probability on the positive ray: as functions of \(\log r\) these normalized logarithms are \(1\)-Lipschitz since powers range from \(0\) to \(N\). At \(r_0\), the pressure changes second derivative. In fact, as \(p\to p_0\),
\[
 I_{h,p}(\alpha)=
 \frac{(1-2h)^2}{2p_0(1-p_0)}(p-p_0)^2+O(|p-p_0|^3),
\]
by the Legendre transform about its mean and \(\Lambda''_{h,p}(0)=p(1-p)\). Reciprocity gives the same nonanalyticity at \(1/r_0\). At \(r=1\) the first derivative with respect to \(\log r\) jumps. At \(p=1/2\), differentiating at the optimizer gives \(\partial_p I_{h,p}(\alpha)=2(1-2h)(1-2\alpha)\). The left slope is \(h+\alpha-2h\alpha<1/2\) and the right slope is its complement.

For any deterministic sequence of centers along which the positive-ray convergence is locally uniform, the subharmonic functions \(u_N=N^{-1}\log|P_L|\) have locally uniform upper bounds, and \(u_N(0)=0\) excludes collapse to minus infinity. Every subsequential \(L^1_{\rm loc}\) limit \(u\), taken in its upper semicontinuous representative, satisfies \(u(z)\le\Phi(|z|)\) for \(z\ne0\) by positivity of coefficients and the submean inequality. Hartogs' lemma gives \(u(r)\ge\limsup u_N(r)=\Phi(r)\) for \(r>0\). If some fixed disk about one of the three points had \(o(N)\) zeros along a subsequence, the distributional Laplacian of a limit would vanish there. Its restriction to the positive ray would then be real analytic, contrary to the respective derivative jump. To obtain a uniform \(\gamma>0\) in probability, choose deterministic tolerances decreasing sufficiently slowly on growing compact intervals so that local uniform convergence holds along every sequence drawn from the resulting sets of probability \(1-o(1)\). A vanishing infimum of the normalized disk count on these sets would give the excluded subsequence.

Finally, put \(r=|z|\), \(p=r/(1+r)<p_0\). The absolute difference between \(P_L(z)\) and \((1+z)^N\) is bounded by \((1+r)^N\Pr[W>(\alpha-\delta)N]\) for every center. The strict displayed inequality makes this exponentially smaller than \(|1+z|^N\), uniformly on compact subsets by the Chernoff bound with sufficiently small \(\delta>0\). Near zero the same bound follows from the deterministic complete initial distance shells through \((\alpha-h-\delta)N\). The identity \(P_L(z)=z^N P_L(1/z)\) gives the reciprocal conclusion.
\end{proof}

The statistic here fixes one accepted center and asks for the probability that its random perturbation is again accepted. Worst-case approximation of partial truth tables by small circuits is studied by \href{https://eccc.weizmann.ac.il/eccc-reports/1995/TR95-041/index.html}{Andreev, Clementi and Rolim}.

\paragraph{Local density at the three positive points.}
Use the local coordinate \(t=\log z\) and put \(f(t)=\Phi(e^t)\) on the real axis. For \(t_0=\log r_0,0,-\log r_0\), let \(Z_N(t_0,\rho)\) count zeros with \(|\log z-t_0|<\rho\), using the branch real on the positive axis. For suitable constants \(b_1,b_2>0\) and every sufficiently small fixed \(\rho>0\),
\[
 \Pr\!\left[b_1\rho^2\le Z_N(t_0,\rho)/N\le b_2\rho^2\right]\longrightarrow1
 \qquad(t_0=\pm\log r_0).
\]
At the central point the exact first-order density is
\[
 \lim_{\rho\downarrow0}\limsup_{N\to\infty}
 \Pr\!\left[\left|\frac{Z_N(0,\rho)}{N\rho}
            -\frac{(1-2h)(1-2\alpha)}{\pi}\right|>\epsilon\right]=0
 \qquad(\epsilon>0).
\]
In the same ordered limits the rescaled zero measures at \(0\), divided by \(N\rho\), converge vaguely in probability to \((1-2h)(1-2\alpha)/(2\pi)\) times length on the imaginary axis.

To see this, take any locally uniform positive-ray sequence from the preceding proof and any limit \(u\) of its normalized potentials in the log coordinate. Then \(u(t)\le f(\Re t)\), with equality on the real axis. These two facts suffice uniformly over all such limits. Write
\[
 v_\rho(w)=\rho^{-j}\big[u(t_0+\rho w)-f(t_0)-\rho b\Re w\big].
\]
Use \(j=2,\ b=f'(t_0)\) at the outer points and \(j=1,\ b=0\) at the center. As \(\rho\downarrow0\) these subharmonic functions are bounded above uniformly on compact sets by Taylor expansion on each side, and \(v_\rho(0)=0\). They therefore have \(L^1_{\rm loc}\) subsequential limits even if the underlying \(u\) varies. Each limit is bounded above by the corresponding piecewise Taylor term evaluated at \(\Re w\), by the submean inequality, and agrees with it on the real axis by Hartogs' lemma. At an outer point the trace consists of two quadratic polynomials with unequal second derivatives. A limit cannot be harmonic on any disk about zero. Consequently the Riesz masses of \(v_\rho\) on the unit disk are bounded away from zero uniformly: otherwise a limiting mass would vanish on its interior. They are bounded above uniformly by Jensen's formula using the upper bound on the disk of radius two and \(v_\rho(0)=0\). Undoing the scaling multiplies these masses by \(\rho^2\).

At the center the upper bound for every limit is \(f'(0-)\Re w\) on the left half-plane and \(f'(0+)\Re w\) on the right. Each affine bound is attained at interior real points of its respective half-plane. The strong maximum principle gives equality there. Thus the tangent potential is unique in \(L^1_{\rm loc}\), and its Riesz measure is the stated line measure because \(f'(0+)-f'(0-)=(1-2h)(1-2\alpha)\). The unit circle has zero mass for this measure. Vague convergence of Riesz measures now transfers these bounds to the finite polynomials in the stated ordered limits. Explicitly, at fixed \(\rho\) compare an open disk to slightly smaller and larger disks of limiting potentials by Portmanteau. For the quadratic orders one may use radii \(\rho/2,2\rho\); for the central constant first use \((1\pm\eta)\rho\) and then send \(\eta\downarrow0\). Compactly supported continuous tests give the vague statement directly. The typical sets of centers used for local uniform convergence have probability tending to one.

\section{Individual distance coefficients on dense center layers}
Keep \(c,\alpha,h,B,L=B\oplus U\) as in the preceding fixed-weight theorem and let \(b_d(L)=[z^d]P_L(z)\). Put
\[
 \beta_0=\frac{\alpha-h}{1-2h},\qquad
 \beta_1=h+(1-2h)\alpha .
\]
For joint binary types with \(\Pr[U=1]=h,\ \Pr[F=1]=q,\ \Pr[F\ne U]=\beta\), the conditional entropy of \(F\) given \(U\) is
\[
 J(q,\beta)=(1-h)H\left(\frac{q+\beta-h}{2(1-h)}\right)
  +hH\left(\frac{q-\beta+h}{2h}\right).
\]
Set \(J=-\infty\) for infeasible types. Define a continuous symmetric function \(\sigma(\beta)=\sigma(1-\beta)\) by
\[
 \sigma(\beta)=
 \begin{cases}
 H(\beta),&0\le\beta\le\beta_0,\\
 J(\alpha,\beta),&\beta_0\le\beta\le\beta_1,\\
 H(\alpha),&\beta_1\le\beta\le1/2 .
 \end{cases}
\]
\begin{theorem}
For every fixed \(\epsilon>0\),
\[
 \Pr_L\!\left[\max_{0\le d\le N}
   \left|N^{-1}\log b_d(L)-\sigma(d/N)\right|>\epsilon\right]\longrightarrow0 ,
\]
with \(\log0=-\infty\). For every center in the layer, the three-branch profile is a uniform lower bound up to \(o(1)\). In particular the plateau \(\beta_1\le d/N\le1-\beta_1\) reaches the full circuit census exponent for every such center.
\end{theorem}
\begin{proof}
By complementation it suffices to take \(d/N\to\beta\le1/2\). For the upper bound work with \(B\oplus\mathcal A\) about \(U\). A fixed word of weight \(qN\) lies at distance \(d\) with probability \(\exp(N[J(q,\beta)-H(q)+o(1)])\) by counting joint types, with \(h_N,d/N\) replacing \(h,\beta\) before taking limits. There are at most \(\exp(N\min(H(q),H(\alpha))+o(N))\) such words in the class. Hence the upper exponential rate for \(\mathbb E b_d\) is at most
\[
 \max_q\{J(q,\beta)-(H(q)-H(\alpha))_+\}.
\]
The function \(J(\cdot,\beta)\) is concave and is maximized at \(q_F=h+(1-2h)\beta\), with value \(H(\beta)\), by independence of \(U\) and the error \(F\oplus U\).
Also \(J(q,\beta)-H(q)=-I(F;U)\) is concave in \(q\), by convexity of mutual information in the channel for a fixed input distribution (the log-sum inequality). If \(\beta\ge h\), it is maximized at \(q_I=(\beta-h)/(1-2h)\), where \(U,F\) are independent. If \(\beta<h\), it is maximized at \(q_I=(h-\beta)/(1-2\beta)\): here \(F\) and the error are independent, attaining the bound \(H(U\mid F)\le H(\beta)\). When \(\beta_0\le\beta\le\beta_1\), these identities give \(q_I\le\alpha\le q_F\le1/2\). Thus the maximum on the middle weight interval \([\alpha,1-\alpha]\) is at \(\alpha\). On \(q\le\alpha\) the conditional entropy is increasing, and on feasible \(q\ge1-\alpha\) it is decreasing. If the latter interval is feasible, its inner endpoint is bounded by the middle-interval comparison. This gives \(J(\alpha,\beta)\) as upper bound. For the first branch use \(J\le H(\beta)\) and for the last use the total census. Markov's inequality bounds the probability of exceeding the resulting rate by any fixed positive tolerance by \(\exp(-\Omega(N))\).

On the first two branches, complete light layers already give the lower bound for every \(U\). Count the words of joint types with \(q<\alpha\) tending to \(\min(q_F,\alpha)\), first leaving fixed slack for synthesis and then shrinking it. Rounding the weight with the required parity changes it by at most two. For \(0<\beta\le\beta_0\) the independent-error type is interior; throughout \([\beta_0,\beta_1]\), \(\alpha\) is a feasible interior weight since \(\beta_0>\alpha-h\) and \(\beta_1<\alpha+h\), with all parameters below \(1/2\). At \(\beta=0\) every modification of \(o(N)\) entries stays below the complete low-weight cutoff for \(B\oplus\mathcal A\).

For the plateau, choose integer weights \(w_N\) within one of
\((d-h_NN)/(1-2h_N)\), so \(w_N/N\to q\in[\alpha,1/2]\).
For any \(\epsilon>0\), at least \(\exp(N[H(\alpha)-\epsilon])\) distinct words of weight \(w_N\) have cost at most \((c-\epsilon')N/n\) for some \(\epsilon'>0\). We give the argument to specify the slice census being used. Partition by \(O(\log n)\) prefix bits into \(M\) dyadic blocks, with \(M\to\infty\) polynomial in \(n\). Fix \(\gamma<\alpha\) close to \(\alpha\). On \(M-1\) blocks use arbitrary words of relative weight \(\gamma\) or \(1-\gamma\), rounding symmetrically. Select the number of heavy blocks so a comparator on the remaining block achieves the exact total \(w_N\): successive totals differ by at most the block length and \(q\) is interior to \((\gamma,1-\gamma)\). Stirling gives \(\exp((H(\gamma)+o(1))N)\) choices. Lupanov's slice synthesis bound \cite{lup,kors} and an \(O(M)\)-gate multiplexer give cost at most \((H(\gamma)/\log2+o(1))N/n\). The classical synthesis theorem provides the block circuits; the weight adjustment and subsequent distance adjustment require the following separate argument.

Reserve a fixed dyadic block of \(R=2^{\lfloor 2n/3\rfloor}\) addresses. Fix any \(U\) in the layer and apply one uniformly random invertible affine permutation of the raw addresses to the selected word family. For each fixed selected word its expected full distance to \(U\) is \(d+O(1)\), and its expected distance on the reserved block lies between \(\min(w_N/N,1-w_N/N)R\) and \(\max(w_N/N,1-w_N/N)R\). These bounds stay a fixed positive fraction of \(R\) away from both endpoints. Its outside distance has variance \(O(N)\): each membership indicator has variance at most one and the covariance of two distinct permuted indicators is \(O(1/N)\) by two-transitivity. Taking complements on prescribed sites changes only covariance signs. Chebyshev gives probability \(1-O(N/R^2)\) that the outside distance belongs to \([d-R,d]\). By averaging there is one permutation, common to at least half the selected words, for which they have outside distance in the desired range. Precomposing an address circuit with this permutation costs \(O(n^2)\) gates. For each such word rewrite the reserved block to give total distance exactly \(d\). This costs \(O(R+n)=o(N/n)\) by a binary decision tree and a prefix switch. Distinct outputs lose a factor of at most \(2^R\) in count. XOR with \(B\) places these modified words in \(\mathcal A\), preserving distance to \(L\). This proves the lower bound at each plateau distance.

For simultaneity, the asymptotic bounds for the mean, the deterministic lower bounds, and the plateau word counts are uniform at fixed accuracy in their respective closed ranges: otherwise compactness gives a convergent subsequence of \(d/N\) contradicting the corresponding type estimates with fixed entropy slack. At either common endpoint both adjacent formulas agree. The lower bounds now hold for every \(U\). Markov and a union bound over \(N+1\) coefficients give the uniform upper with exponentially small exceptional fraction at each fixed tolerance.
\end{proof}

\section{A reference chosen from all truth tables}
The preceding fixed-positive-weight result used a typical \emph{accepted} reference. We now allow $L$ to be uniformly random among all $2^N$ tables, at a fixed gate rate $sn/N\to c\in(0,1)$. Such a table is usually rejected. The enumerator is still $P_L(z)=\sum_{F\in\mathcal A}z^{d(F,L)}$. Its zero multiplicity at the origin is the distance of $L$ to the easy class. Put $\delta=H_2^{-1}(1-c)$ on the increasing branch and $r_-=\delta/(1-\delta)$.

\begin{theorem}
Let $k/N\to\beta$ be an integer distance sequence with $H_2(\beta)>1-c$. At this distance every reference $L$, of any complexity, has at least $2^{(c+H_2(\beta)-1-o(1))N}$ easy neighbors. For uniform $L$ their number has this logarithmic rate in probability. If $H_2(\beta)<1-c$, the probability that a uniform $L$ has any such neighbor is $2^{-(1-c-H_2(\beta)+o(1))N}$. The last assertion includes $\beta=0,1$. The covering radius of $\mathcal A$ and the distance of a uniform reference to $\mathcal A$ are $(\delta+o(1))N$, the latter in probability.
\end{theorem}
\begin{proof}
For the universal lower bound, prescribe $k$ changes to $L$. Round so that the fractions of changed zeros and changed ones differ from $k/N$ only by integer rounding in their respective class counts. The resulting joint distribution has conditional output entropy $H_2(\beta)+o(1)$, since there are only two source classes and average rounding variation is $O(1/N)$. Its mutual information is $H(Y)-H_2(\beta)+o(1)\le1-H_2(\beta)+o(1)$. The constructive bound in the companion manuscript \emph{How many small Boolean circuits realize a prescribed joint type?}, Theorem~2, supplies the claimed list. Indeed for any fixed exponent $e<c+H_2(\beta)-1$, spend an additional normalized rate $e$ above that mutual information; it is at most the available conditional entropy. Taking $e$ up to the stated bound is uniform over references. This covering-radius upper is also the classical scalar result of Pippenger \cite{pipp}.

For every $F$, a uniform $L$ has distance $k$ with probability $2^{-N}\binom Nk$. Multiplication by $|\mathcal A|\le2^{(c+o(1))N}$ and Markov give the typical upper above threshold and the rare-event upper below. For the matching lower below threshold, first condition $L$ to have $\lfloor N/2\rfloor$ ones, an event of probability $2^{-o(N)}$. Prescribe $k$ changes nearly equally on its two classes. Then $H(Y)\to1$, $H(Y\mid L)\to H_2(\beta)$, and the below-threshold part of the companion Theorem~1 gives the claimed probability from below. When $\beta$ equals zero or one the same limiting argument applies. Union bounding the elementary expectation over all $k\le(\delta-\epsilon)N$ excludes easy neighbors with high probability there, for each fixed positive $\epsilon<\delta$. This and the universal lower on a layer a little above $\delta N$ give the distance assertions.
\end{proof}

\paragraph{A positive-axis formula and two zero accumulations.}
In the log coordinate $t$ on the real axis, the limiting pressure $f(t)=\lim N^{-1}\log P_L(e^t)$ for a uniform reference, in probability, is
\[
 f(t)=
 \begin{cases}
   \delta t, &t\le\log r_-,\\
   (c-1)\log2+\log(1+e^t), &\log r_-\le t\le-\log r_-,\\
   (1-\delta)t, &t\ge-\log r_- .
 \end{cases}
\]
Indeed Markov's bound on each shell holds simultaneously on all shells with a fixed positive exponent slack, by a union bound over $N+1$ terms. It excludes macroscopically negative-exponent shells. The uniform lower on each strictly feasible layer then gives, by taking nearby layers at the endpoints, the maximum of
$(c-1)\log2+H_2(\beta)\log2+\beta t$ over $\delta\le\beta\le1-\delta$. Maximization in $\beta$ gives the formula. Lipschitz continuity extends convergence locally uniformly on the real axis in probability.

With probability tending to one, disks of each sufficiently small fixed radius $\rho$ about $t=\pm\log r_-$ contain $\Theta(\rho^2 N)$ zeros, in the log coordinate. The constants are uniform for sufficiently small $\rho$, taking the large-$N$ limit first. This follows by the same subharmonic rescaling argument as at the two noise transitions in the fixed-weight analysis: the two one-sided Taylor quadratics have unequal second derivatives. On either open half-plane $\Re t<\log r_-$ or $\Re t>-\log r_-$ the radial upper envelope is harmonic and every limit attains it on the real axis. Thus the strong maximum principle identifies the potential there. In particular any fixed compact subannulus of $0<|z|<r_-$ or $1/r_-<|z|<\infty$ contains $o(N)$ zeros in probability. At $z=0$ itself the multiplicity is $(\delta+o(1))N$. The real-axis formula gives no zero accumulation requirement at $z=1$ for this choice of reference.

\section{Centering at a uniformly chosen accepted table}
The low-complexity centers of Theorem~1 have a common local zero law. We now ask what happens at a center sampled uniformly from all distinct accepted truth tables. Write $R_L$ for the least modulus of a zero of $P_L$.
\begin{theorem}\label{th:typical}
Suppose $n=o(s)$ and $s=o(N/n)$, and choose $L$ uniformly in $\mathcal A$. For each positive $r=o(1/\sqrt{NK})$,
\[
 \Pr[R_L>r]\longrightarrow 1,\qquad P_L(r)\longrightarrow 1
 \quad\text{in probability}.
\]
If $K^3=o(N)$, then $R_L/(K/N)\to\infty$ in probability. In the same regime, at flip probability $p=xK/N$ for fixed $x>0$, the acceptance probability conditioned on $L$, divided by $(1-p)^N$, tends to one in probability.
\end{theorem}
\begin{proof}
We use the edge-isoperimetric inequalities for powers of the cube proved by Rashtchian and Raynaud \cite[Theorems~2 and~3]{rr}. In a general family $A$ of $M$ strings, let $B=\log_2 M=o(N)$ and $b=\min(\lceil2B/\log_2(N/B)\rceil,\lfloor B\rfloor)$. For $1\le t\le B$, the average number of other family members at distance $2t$ from a uniform member is at most $2(8e/t)^{2t}(Nb)^t$. At distance $2t+1$, an upper bound is $2(16e/(2t+1))^{2t+1}(Nb)^t B$. These results bound all pairs at distance up to the specified radius, hence bound each individual distance. The usual cube edge inequality bounds the average distance-one count by $B$.

Write $u=r\sqrt{Nb}$. If $u=o(1)$, summation of the even terms through $2\lfloor B\rfloor$ contributes $O(u^2)$ to $\mathbb E_L(P_L(r)-1)$. The odd terms through $2\lfloor B\rfloor+1$ contribute $O(rB)$. The remaining terms sum to at most $M r^{2\lfloor B\rfloor+2}$. In our application the complete early shell and the circuit census imply $B\sim K\log_2(N/K)$, so $b\sim 2K$, $B=o(\sqrt{NK})$, and these bounds tend to zero. Markov's inequality gives $P_L(r)-1=o(1)$ in probability. Whenever $P_L(r)-1<1$, the triangle inequality excludes a zero on $|z|\le r$, since $P_L(0)=1$. Finally $K/N=o(1/\sqrt{NK})$ under the stated cubic condition. Apply the first assertion at any fixed multiple of $K/N$, including the flip odds $p/(1-p)$.
\end{proof}

The radius scale is sharp among general families with the same entropy. For the family of weight-$k$ strings with $k\to\infty$, $k=o(N)$, every center has polynomial $\sum_j\binom{k}{j}\binom{N-k}{j}z^{2j}$. With $z=w/\sqrt{k(N-k)}$ it converges locally uniformly to $\sum_{j\ge0}w^{2j}/(j!)^2=I_0(2w)$, by the same summable majorant on each compact disk. This limit has a zero at $w=i j_{0,1}/2$, where $j_{0,1}$ is the first positive zero of $J_0$. Local uniform convergence supplies zeros of the displayed order. Including the complementary sphere gives the same limit, because its contribution on compact disks is bounded by $\binom Nk |z|^{N-2k}=o(1)$. Determining the nearest-zero order at a typical MCSP center beyond our one-sided bound remains open.

\section{Centering at a rigid easy table}
For $n\le s$ and $s\log(s+2)=o(\sqrt N)$, the appendix constructs an exact-size-$s$ table $L$ at distance at least $\Delta=2^{m-3}$ from any other member of $\mathcal A$, where $m=n-\log_2(s\log_2(s+2))+O(1)$. Thus for $|z|<1$,
\[
 |P_L(z)-1|\le (|\mathcal A|-1)|z|^\Delta.
\]
Complement symmetry gives $P_L(z)=z^N P_L(1/z)$. Every zero therefore lies in the annulus with radii $\exp(\pm \log(|\mathcal A|-1)/\Delta)$. The rooted count makes the exponent $O((s\log s)^2/N)=o(1)$. The angular discrepancy is $O(\sqrt{\log|\mathcal A|/N})$ by the Erd\H{o}s--Tur\'an inequality \cite{et}, applied to the monic polynomial with constant coefficient one. This gives both uniform radial concentration and Hausdorff convergence. Since $K=o(N)$, substituting $z=p/(1-p)=O(K/N)$ in the displayed bound makes the right-hand side tend to zero throughout this rigid-center regime.
\section{Comparison with previous results}
Pritsker and Varga \cite{pv} analyze the Szeg\H{o} limit measure for the normalized sections of the exponential. Janson and Norfolk \cite{jn} analyze zeros of binomial sections, including the location of the negative extremum. The enumerator here has degree $N$ even in the sparse scaling; beyond the complete weight layers, its coefficients can differ from those of any single binomial section. The entropy cap controls that unknown tail, and positivity on the real axis determines the exterior potential. In the fixed-fraction case the second complete tail produces a reciprocal copy of the classical curve and the remaining mass on the unit circle.
\appendix
\section{Keeping the complexity of a hard core}
Let \(y\in\{0,1\}^r,\ z\in\{0,1\}^m,\ r+m=n,\ m\ge2\), and take an arbitrary Boolean \(h(y)\). Set
\[
                   L_h(y,z)=h(y)\oplus\bigoplus_{i=1}^m z_i .
\]
If \(h\) is nonconstant, \(C(L_h)=C(h)+m\). The upper bound appends \(m\) gates. For the reverse inequality fix the \(z\) bits successively: each is essential, and while it remains, the nonconstant \(h\) ensures that the function to which it belongs cannot be a bare literal. The final restriction computes \(h\) up to a constant output sign.

\textbf{Rigidity with the core charged.} If \(1\le k=|\{f\ne L_h\}|\le 2^{m-2}\), then
\[
                    C(f)\ge C(h)+2m-\lfloor\log_2 k\rfloor-3. \tag{A.1}
\]
This is a special gate-elimination statement about a parity extension and a sparse defect. It does not use general additivity for disjoint circuits.

First, while the remaining errors have more than one \(z\)-projection, choose a \(z\) coordinate on which they vary and retain its smaller, nonempty error side. After \(q\le\lfloor\log_2 k\rfloor\) restrictions all errors have the same projection on the remaining \(m'=m-q\ge2\) parity coordinates; the \(y\) values may vary. All these restrictions save at least one gate each. They keep the error fraction below one half, making every remaining parity variable essential, and before each restriction there are at least two such variables so the current function cannot be a literal. Absorb the constant sign change from fixed parity bits into \(h\).

We prove next that a circuit for
\[
 h(y)\oplus {\rm Par}(z)\oplus {\bf1}[y\in B,\ z=u]
\]
on \(m'\) remaining parity coordinates, fixed \(u\) and nonempty \(B\subseteq\{0,1\}^r\), has size at least \(C(h)+2m'-3\). During the reduction keep every signal denoted below as a tag as a separate auxiliary output; its subcircuit has no syntactic path from any remaining \(z_i\). Pruning gates irrelevant to both the main output and all tags is permitted. Unary gates and signs may be absorbed. For each original parity index \(i\), maintain the multiset \(A_i\) of its tags. An additive constant sign may change at any step. With \(I\) the remaining coordinates, the invariant is that the computed main function is
\[
 h(y)\ \oplus\! \bigoplus_{i\in I}z_i\
 \oplus\! \bigoplus_i\bigoplus_{H\in A_i}H(y)
 \ \oplus\ {\bf1}\!\left[y\in B',\
    z_i=u_i\oplus\bigoplus_{H\in A_i}H(y)\ {\rm for\ every}\ i\in I\right]
 \ \oplus c ,
\]
where \(B'\ne\varnothing\). Initially all multisets are empty. The index in the double sum outside the indicator ranges over \emph{all} original parity positions, including removed ones.

While at least three parity variables remain, if one has two or more gate readers, fix it to its value at a chosen point of the nonempty error graph. At least two gates disappear. The baseline changes by a fixed sign and the graph loses that coordinate, narrowing \(B'\) while remaining nonempty. No tags are lost. A parity input cannot be an auxiliary output, and the main output cannot be a bare parity input, since some \(y\)-fiber is a full parity with a singleton defect.

Otherwise every remaining \(z_i\) has just one reader. Take a topologically earliest gate syntactically depending on a \(z_i\). It cannot consume two raw \(z\)'s. At a fixed \(y\in B'\), the computed function of all \(z\)'s is full parity up to sign with one exceptional point. Setting the other parity bits off its projection forces the putative two-private-input gate to be XOR/XNOR; setting them at that projection violates its resulting pair-flip invariance. Thus the earliest gate is \(G(z_i,H(y))\) with a syntactically \(Y\)-only partner. Any nonparity binary gate has a canalizing partner value at which the output is independent of \(z_i\). If \(H\) attains that value, then at that \(y\) the computed function must be independent of \(z_i\) for all \(z\), impossible since this fiber has at most one parity error. If \(H\) never attains the value, it is semantically constant and the gate can be bypassed by a signed raw wire. In the latter case retain a constant sign as a tag if needed.

For a parity gate, perform the shear \(z'_i=G=z_i\oplus H(y)\oplus c_i\), replace the gate's wire everywhere by the new raw \(z'_i\), and retain its signed partner as one new tag in \(A_i\). The singleton graph and baseline change exactly as in the invariant. This removes one gate; retaining the partner as a second output incurs none. No existing tag changes its \(Y\)-only syntax, since it had no \(z_i\)-dependent ancestor in the first place and the rewrite affects only descendants of that raw variable. Each such shear strictly reduces total gate count, so it cannot prevent termination. The constant-partner bypass has the same accounting.

Stop with two remaining parity variables. Fix one of them to zero before reconnecting the other. The fiber at any \(y\in B'\) is a two-variable parity with one defect, so both parity variables have a gate reader and no auxiliary output depends on them. This restriction saves one additional gate. Choose the remaining index \(j\). Feed \(z_j=u_j\oplus1\oplus\bigoplus_{H\in A_j}H(y)\). The graph indicator vanishes for every \(y\). Constructing this replacement signal uses at most \(|A_j|\) additional gates. Topologically it is legal to compute all retained \(Y\)-only tag subcircuits first, then compute the replacement signal, then the remaining circuit gates: no tag has any syntactic dependence on a parity raw wire. In the output, the replacement cancels all tags from \(A_j\). Recover \(h\) up to a constant sign by taking XOR with the tags at all \emph{other} original indices, including indices fixed during the reduction, using at most \(\sum_{i\ne j}|A_i|\) gates. The total rebuilding cost is thus at most the number of single-gate shears and bypasses. The \(m'-2\) multi-reader restrictions saved at least \(2(m'-2)\) further gates. We have constructed \(h\) with at least \(2m'-3\) fewer than the original gates. Adding back the \(q\) first restrictions proves (A.1).
\section{Hitting every integer circuit complexity}
For any positive integer \(a\), let \(B(a)\) denote the maximum full-binary complexity of a Boolean function on \(a\) inputs. On \(a+2\) inputs every integer \(0\le u\le B(a)\) is a circuit complexity. The two extra inputs provide a finite-state computation whose transitions each increase \emph{minimum} circuit complexity by at most one.

Take a function \(G(y)\) of complexity \(B(a)\) and a DNF for it. A straight-line program scans this DNF with two one-bit state registers \(p,q\). Initialize \(p=0\). For each term set \(q=1\), then update \(q\gets q\wedge\ell(y)\) once per literal, and finally \(p\gets p\vee q\). The result is \(G\). For each suffix of this program, regard its final \(p\) value as a function of its starting inputs \(y,p,q\). The empty suffix computes \(p\), of size zero. Prepending one instruction substitutes either a constant or the output of a single binary gate for one formal register, leaving all other raw inputs as they were. A circuit for the old suffix can be used once with that one shared substitution, so the new suffix complexity is at most the old one plus one. The full suffix at initial \(p=0,\ q=1\) gives \(G(y)\), hence its complexity is at least \(B(a)\). From zero to this value, upward steps of at most one necessarily attain each intermediate integer. The DNF length is immaterial.

Elementary circuit counting gives \(B(a)\ge c2^a/a\) for an absolute \(c>0\). Thus for each budget \(s\ge n\), choose \(a=\lceil\log_2(s\log_2(s+2))\rceil+O(1)\) with the constant large enough that \(B(a)>s\); put \(r=a+2,\ m=n-r\). If \(s=o(N/n)\), then \(s\log s=o(N)\), so \(m\to\infty\). Since \(s-m\ge r\ge1\), interpolation supplies a nonconstant \(r\)-bit core \(h\) of complexity exactly \(s-m\). Both \(L_h\) and its complement have complexity exactly \(s\). Finding this core in polynomial time is not part of the existence claim.

For the selected core, (A.1) implies distance at least $2^{m-3}$ to every distinct table of size at most $s$ whose distance is at most $2^{m-2}$. Larger distances satisfy the same conclusion immediately.
\begin{thebibliography}{9}
\bibitem{pipp} N. Pippenger, \href{https://doi.org/10.1007/BF01683269}{Information theory and the complexity of Boolean functions}, Mathematical Systems Theory 10 (1976/77), 129--167.
\bibitem{red} N. P. Red'kin. \href{https://www.mathnet.ru/eng/dm172}{On the complexity of Boolean functions with a small number of ones}, Discrete Mathematics and Applications, 2004.
\bibitem{lup} O. B. Lupanov, On a certain approach to the synthesis of control systems: the principle of local coding, Problemy Kibernetiki 14 (1965), 31--110.
\bibitem{kors} A. D. Korshunov, \href{https://www.mathnet.ru/links/7cc2cdffddc9c901d3badd138b575bb5/rm9459_eng.pdf}{Computational complexity of Boolean functions}, Russian Mathematical Surveys 67 (2012), \S2.4, item 5.
\bibitem{chash} A. V. Chashkin, \href{https://keldysh.ru/papers/2024/mvk/mvk2024_152.pdf}{Mathematical Questions of Cybernetics (2024)}, Theorem~6.1, pp.~199--200.
\bibitem{pv} I. E. Pritsker and R. S. Varga, \href{https://math.okstate.edu/people/igor/96a.pdf}{The Szeg\H{o} curve, zero distribution and weighted approximation}, Trans. Amer. Math. Soc. 349 (1997), 4085--4105.
\bibitem{jn} S. Janson and T. S. Norfolk, \href{https://arxiv.org/abs/0903.3866}{Zeros of sections of the binomial expansion}, Electronic Transactions on Numerical Analysis 36 (2009), 27--38.
\bibitem{rr} C. Rashtchian and W. Raynaud, \href{https://arxiv.org/abs/1909.10435}{Edge isoperimetric inequalities for powers of the hypercube}, Electron. J. Combin. 29 (2022), P1.47.
\bibitem{et} P. Erd\H{o}s and P. Tur\'an, \href{https://renyi.hu/~p_erdos/1950-08.pdf}{On the distribution of roots of polynomials}, Ann. of Math. 51 (1950), 105--119.
\end{thebibliography}
\end{document}
